\section{Role of Human Oversight (Human-in-the-Loop)}

The foundations chapter distinguished human-in-the-loop as a deliberate design choice --- a decision about where, in an otherwise automated process, human judgment is required before the process can continue --- from oversight understood loosely as "a person eventually looked at the result." It also distinguished several forms this can take: automatic acceptance with no check at all, review after a finished output already exists, and a genuine gate that pauses the process and waits for explicit approval before it proceeds. The two parts of the practical project sit at different points along this range, and the difference in where oversight was placed, not merely whether it existed, is what separated the two.

Part 1 had a standing rule requiring the agent to ask before modifying any existing file, but in practice this amounted to review after the fact rather than a gate: changes were generally accepted once they worked, without being checked against the system they were being added to, and oversight only became substantive once a problem was already obvious enough to notice on its own --- by which point correcting it meant undoing an accumulation rather than preventing one decision. \textbf{The practical cost of this placement is not that oversight was absent, but that it arrived too late to be cheap: the same logic that makes a requirements-stage mistake far less costly to fix than one caught after implementation applies directly here, and review positioned after output already exists forfeits exactly that advantage.}

Part 2 moved the same kind of judgment earlier, into a genuine gate rather than a later check: a human is required to approve the Design Review portion of a plan before any implementation begins, and is explicitly instructed to reject a plan outright --- not correct it after the fact --- if the reasoning behind an architectural choice is vague or if it is unclear which component or service owns a given piece of state or logic. \textbf{The recommendation this supports is to place mandatory approval at the point where a decision is still cheap to reverse --- before code exists to be reworked --- rather than relying on review to catch the same problem once it has already been built on top of.}

The division of labor also shifted, and this shift is itself a form of oversight design worth naming as a practice. In Part 1, a human did the work of noticing that something was wrong in the first place, and the agent's role was limited to executing a correction once told. In Part 2, the Architect Agent surfaced at least one significant concern --- flagged inside a plan, before any code existed --- that a human then reviewed and approved rather than discovered independently, and the Post-Implementation Review Agent's report narrowed the human's role after implementation to deciding whether to fix, defer, or accept a finding the agent had already produced on its own. \textbf{This suggests that human authority is best reserved for judgment a structurally different party cannot supply --- approving or rejecting a plan, weighing a trade-off, deciding what to do about a flagged problem --- while the work of detecting candidate problems in the first place can be delegated to an agent role built specifically for that check, provided a human still holds the authority to act on what it finds.}

